\section{导数的运算规则}
\label{sec:ru-derivative-rules}
\label{sec:analysis-scalar-chain-rules}

直接用定义求导，每个函数都要重新计算差商。若已知几个函数的导数，由它们组合出的函数的导数应能直接算出。
本节在导数已经存在的条件下，分别研究四则运算、复合和取反函数怎样传递到导数，下一节再用这些规则计算基本函数的导数。
以下设$f,g$在所讨论的点$a$的某个开邻域内有定义，并在$a$处可导。

\subsection{导数的四则运算}
\label{sec:ru-derivative-arithmetic}

和的差商等于两个差商之和，常数倍的差商等于差商的同一常数倍，令$h\to0$即得$(f\pm g)'(a)=f'(a)\pm g'(a)$与$(cf)'(a)=cf'(a)$。
乘积的增量可以拆成三部分。记$\Delta f=f(a+h)-f(a)$、$\Delta g=g(a+h)-g(a)$，则
\[
 f(a+h)g(a+h)-f(a)g(a)=g(a)\,\Delta f+f(a)\,\Delta g+\Delta f\,\Delta g.
\]
由一次展开，前两项之和为$[f'(a)g(a)+f(a)g'(a)]h+o(h)$；两个增量都至多与$|h|$同阶，乘积$\Delta f\,\Delta g$与$h^2$同阶，所以也是$o(h)$。读出一次项的系数，得到乘积法则
\begin{equation}
 (fg)'(a)=f'(a)g(a)+f(a)g'(a).
 \label{eq:analysis-scalar-sum-product-rules}
\end{equation}
增量的乘积之所以可以并入余项，是因为两个小量相乘比每个小量都小得多；这一步使乘积的导数只剩两项。

商需要分母非零。设$g(a)\neq0$，可导蕴含连续，所以$g$在$a$附近不为零。先求倒数的导数：
\[
 \frac1h\Bigl(\frac1{g(a+h)}-\frac1{g(a)}\Bigr)=-\frac{g(a+h)-g(a)}h\cdot\frac1{g(a+h)g(a)}\to-\frac{g'(a)}{g(a)^2},
\]
再把$f/g$写成$f\cdot(1/g)$用乘积法则，得到
\[
 \Bigl(\frac fg\Bigr)'(a)=\frac{f'(a)g(a)-f(a)g'(a)}{g(a)^2}.
\]

\subsection{复合函数的求导法则}
\label{sec:ru-chain-rule}

复合函数的输入变化先经过内层、再传入外层，总的一次变化由两层的变化率相乘得到。
\begin{theorem}[链式法则]
\label{thm:ru-chain-rule}
若$f$在$a$处可导，$g$在$f(a)$处可导，且$g\circ f$在$a$的某个邻域内有定义，则
\begin{equation}
 (g\circ f)'(a)=g'(f(a))\,f'(a).
 \label{eq:analysis-scalar-chain-rule}
\end{equation}
\end{theorem}
\begin{proof}
记$b=f(a)$、$k(h)=f(a+h)-f(a)$。由$f$的一次展开，$k(h)=f'(a)h+o(h)$，所以$k(h)\to0$，并且充分小的$h$满足$|k(h)|\leq(|f'(a)|+1)|h|$。
把$g$在$b$处的一次展开写成$g(b+k)=g(b)+g'(b)k+|k|\,\varepsilon(k)$，其中$k\neq0$时$\varepsilon(k)=[g(b+k)-g(b)-g'(b)k]/|k|\to0$，并规定$\varepsilon(0)=0$。代入$k=k(h)$，
\[
 g(f(a+h))-g(f(a))=g'(b)f'(a)h+g'(b)\,o(h)+|k(h)|\,\varepsilon(k(h)).
\]
最后一项的绝对值不超过$(|f'(a)|+1)|h|\,|\varepsilon(k(h))|$，也是$o(h)$。所以一次项系数为$g'(b)f'(a)$。
\end{proof}
证明中规定$\varepsilon(0)=0$，是为了处理内层增量$k(h)$可能为零的情形；若照搬记号$\frac{\mathrm dg}{\mathrm dx}=\frac{\mathrm dg}{\mathrm dy}\cdot\frac{\mathrm dy}{\mathrm dx}$“约去$\mathrm dy$”，就会在$k(h)=0$时除以零。
记号可以帮助记忆，但不能代替证明。用一个能直接展开的例子核对：$(1+x^2)^3$以$y=1+x^2$为中间变量，链式法则给出$3y^2\cdot2x=6x(1+x^2)^2$；
直接展开$(1+x^2)^3=1+3x^2+3x^4+x^6$再求导，得$6x+12x^3+6x^5$，二者相同。

\subsection{反函数的求导法则}
\label{sec:ru-inverse-derivative}

上一章已经证明，区间上连续且严格单调的函数有连续的反函数（定理~\ref{thm:ru-inverse-continuity}）。现在判断反函数何时可导。
设$f$在区间$I$上连续、严格单调，$g=f^{-1}$。设$a$是$I$的内点，$f$在$a$处可导且$f'(a)\neq0$，记$b=f(a)$，$b$是区间$f(I)$的内点。
对充分小的$k\neq0$，令$h=g(b+k)-a$。$g$是单射，所以$h\neq0$；$g$连续，所以$k\to0$时$h\to0$。由$f(a+h)=b+k$，
\[
 \frac{g(b+k)-g(b)}k=\frac h{f(a+h)-f(a)}\to\frac1{f'(a)}.
\]
因此$g$在$b$处可导，且$g'(b)=1/f'(a)$，即$(f^{-1})'(b)=1/f'(f^{-1}(b))$。

对$f(g(y))=y$用链式法则也能得到同一公式，但那要求事先知道$g$可导，所以上面先用差商证明可导性；证明中$h\to0$这一步依赖反函数的连续性。
条件$f'(a)\neq0$不能省略：$x^3$在$0$处导数为零，反函数$\sqrt[3]y$在$0$处的差商为$1/\sqrt[3]{k^2}$，$k\to0$时无界，反函数不可导，图像在原点有一条竖直的切线。
在区间端点只能取单侧差商，结论相应地改为单侧导数。下一节用这一规则求对数与反三角函数的导数。
